% SGA 3，第 VII_B 篇，英文重建稿——第 1 轮。
% 范围：权威底本分篇第 40--41 页，命题 1.6.E 全文，止于编者注 73 及命题 1.6.F 之前。
% 合并阅读本页码： 553--554.
% 再版印刷页码： 542--543.
% 本范围内原版页边定位： 511.
% 权威底本： Polo--Gille Exp7B-13oct24.pdf,
% SHA-256 D303690E0195654BBE7C3DAA5D89EA23E24CD709A3575217234AC19AE8836694.
% 曾参考 Jacob Reinhold 的 CC BY 4.0 Markdown 作为比对谱系。
% 本范围内没有图表；全部数学均为原生 LaTeX。
\medskip\noindent\textbf{命题 1.6.E。}\phantomsection \label{sga3:VIIB:1.6.E}\textit{(a) 设 \(X\) 为 \(k\) 上的形式簇，仿射代数为 \(A\)。下列条件等价：\begin{enumerate}\item[\textup{(i)}]\(X\) 在 \(k\) 上 平展。\item[\textup{\SGAInlineRefTarget{sga3:VIIB:1.6.E:item:ii}(ii)}]\(X\) 在 \(k\) 上拓扑平坦，且对角态射 \[
    \Delta:X\longrightarrow X\times X
  \] 为局部同构；即对 \(X\) 的每个点 \(x\)，\(\Delta_X\) 诱导同构 \[
    \mathcal O_{X\times X,\Delta(x)}
    \xrightarrow{\ \sim\ }
    \mathcal O_{X,x}.
  \]。\item[\textup{\SGAInlineRefTarget{sga3:VIIB:1.6.E:item:iii}(iii)}]对每个投影到 \(s\in\operatorname{Spf}(k)\) 的 \(x\in X\)，存在同构 \[
    \mathcal O_{X,x}\simeq k_s[T]/(F),
  \]，其中 \(F\in k_s[T]\) 为可分首一多项式 \textup{（参见 \SGARefLink{sga3:VIIB:1.6.A}{1.6.A}~\SGARefLink{sga3:VIIB:1.6.E:item:ii}{(ii)}）}。\item[\textup{\SGAInlineRefTarget{sga3:VIIB:1.6.E:item:iv}(iv)}]\(X\) 在 \(k\) 上拓扑平坦，且对每个投影到 \(s\in\operatorname{Spf}(k)\) 的 \(x\in X\)，\[
    \mathcal O_{X,x}\widehat\otimes_k\kappa(s)
  \] 为 \(\kappa(s)\) 的有限可分扩张。\item[\textup{\SGAInlineRefTarget{sga3:VIIB:1.6.E:item:v}(v)}]对 \(k\) 的每个开理想 \(\mathfrak l\)，\[
    A\widehat\otimes_k(k/\mathfrak l)
  \] 的每个局部分量均为 Artin 环 \(k/\mathfrak l\) 上的有限 平展 代数。\end{enumerate}(b) 范畴 \(\operatorname{Vaf}^{\mathrm{\acute{e}t}}/\kappa(k)\) 可与 \(S_{\kappa(k)}\) 上的 平展 概形范畴等同 \textup{（参见 \SGARefLink{sga3:VIIB:1.6.C}{1.6.C}）}，且函子 \(X\mapsto X_\kappa\) 诱导等价 \[
  \operatorname{Vaf}^{\mathrm{\acute{e}t}}/k
  \xrightarrow{\ \sim\ }
  \operatorname{Vaf}^{\mathrm{\acute{e}t}}/\kappa(k)
\] \textup{（参见 SGA~1，\SGARefLink{sga3:i:definition:6-1}{I 6.1}）}。}

\medskip\noindent\textit{证明。}设 \(I\) 为乘法态射 \[
  m:A\widehat\otimes_kA\longrightarrow A.
\] 的核。先假设 \(X\) 在 \(k\) 上 平展，故 \(A\) 在 \(k\) 上形式 平展。由 EGA \(0_{\mathrm{IV}}\)，20.7.4，\(I/I^2\) 相对于由 \(I\) 诱导的商拓扑的分离完备化为零；故 \[
  I=I^2.
\]。对每个 \(x\in X\)，理想 \(I\) 包含于 \(A\widehat\otimes_kA\) 的极大理想 \(\mathfrak m_{\Delta(x)}\)。故其局部化 \(I_{\mathfrak m_{\Delta(x)}}\) 包含于 \(\mathcal O_{X\times X,\Delta(x)}\) 的极大理想。于是 Nakayama 的 \SGARefLink{sga3:VIIB:0.3}{引理 0.3} 给出 \[
  I_{\mathfrak m_{\Delta(x)}}=0.
\]。因此 \(\Delta:X\to X\times X\) 为局部同构。

其次假设 \(\Delta\) 为局部同构，且 \(k\) 是一个域 \(\kappa\)。我们证明每个 \(\mathcal O_{X,x}\) 均为有限维 平展 \(\kappa\)-代数。以 \(\operatorname{Spf}(\mathcal O_{X,x})\) 代替 \(X\)，可以假设 \[
  A=\mathcal O_{X,x}
\] 为局部环。我们按照 EGA \(\mathrm{IV}_4\)，17.4.1，\textup{(b)}\(\Rightarrow\)\textup{(d'')} 的证明进行。设 \(K\) 为包含剩余域 \(\kappa(x)\) 的 \(\kappa\) 的有限正规扩张，置 \[
  B=A\otimes_\kappa K
   =A\widehat\otimes_\kappa K,
  \qquad
  X_K=\operatorname{Spf}(B)
      =X\widehat\otimes_\kappa K.
\]。由于 \([K:\kappa]<\infty\)，这两个张量积相同。

态射 \[
  \Delta_K:
  X_K\longrightarrow X_K\times_KX_K
\] 仍为局部同构。因此对每个 \(y\in X_K\)，乘法诱导同构 \[
  \left(
    B_y\widehat\otimes_KB_y
  \right)_{\mathfrak m_{\Delta_K(y)}}
  \xrightarrow{\ \sim\ }
  B_y.
\]。\(B_y\) 的剩余域为 \(K\)；例如参见 \(\mathrm{VI}_A\)，\SGARefLink{sga3:VIIA:1.1.1}{1.1.1}，编者注~(11)。此外，\[
  C=B_y\widehat\otimes_KB_y
\] 本身已经是局部环。事实上，\[
  \mathfrak n
  =
  \mathfrak m_y\widehat\otimes_KB_y
  +
  B_y\widehat\otimes_K\mathfrak m_y
\] 由拓扑幂零元构成，因而包含于 \(C\) 的根中，而 \[
  C/\mathfrak n
  =
  K\widehat\otimes_KK
  =
  K
\] 是域。因此乘法本身为同构 \[
  B_y\widehat\otimes_KB_y
  \xrightarrow{\ \sim\ }
  B_y.
\]。选取 \(B_y\) 在 \(K\) 上包含单位元 \(1\) 的伪基，便得 \(B_y=K\)。由于 \(B\) 在 \(A\) 上有限，集合 \(X_K\) 有限。因此 \[
  B=A\otimes_\kappa K
\] 是有限多个 \(K\) 的副本的乘积，且 \(A\) 为有限 平展 \(\kappa\)-代数。

这样已证明，在域 \(\kappa\) 上，每个在 \(\kappa\) 上 平展 的预有限 \(\kappa\)-代数 \(A\) 都是有限可分扩张 \(K_i/\kappa\) 的乘积，赋予乘积拓扑。因此 \(\operatorname{Spf}(A)\) 为各概形 \[
  \operatorname{Spf}(K_i)=\operatorname{Spec}(K_i),
\] 的直和，且 \(\operatorname{Vaf}^{\mathrm{\acute{e}t}}/\kappa\) 可与 平展 \(\kappa\)-概形范畴等同。

上述论证证明了 \textup{\SGARefLink{sga3:VIIB:1.6.E:item:ii}{(ii)}}\(\Rightarrow\)\textup{\SGARefLink{sga3:VIIB:1.6.E:item:iv}{(iv)}}，因为 \(\mathcal O_{X,x}\widehat\otimes_k\kappa(s)\) 在 \(\kappa(s)\) 上形式 平展，同时也证明了 \textup{(b)}。事实上，回到任意伪紧环 \(k\) 的情形。刚才已将范畴 \(\operatorname{Vaf}^{\mathrm{\acute{e}t}}/\kappa(k)\) 与 \(S_{\kappa(k)}\) 上的 平展 概形范畴等同。由 \SGARefLink{sga3:VIIB:1.6.D}{1.6.D}，为证明 \[
  \operatorname{Vaf}^{\mathrm{\acute{e}t}}/k
  \longrightarrow
  \operatorname{Vaf}^{\mathrm{\acute{e}t}}/\kappa(k)
\] 是等价，只需证明本质满射性。

设 \(s\in\operatorname{Spf}(k)\)，且 \(K/\kappa(s)\) 为有限可分扩张。选取本原元 \(\xi\in K\)，令 \(n=[K:\kappa(s)]\)，并选取次数为 \(n\) 的首一多项式 \[
  F\in k_s[T]
\]，其像 \(\overline F\in\kappa(s)[T]\) 为 \(\xi\) 的极小多项式。由于 \(\overline F'\) 在 \(\kappa(s)[T]/(\overline F)\) 中可逆，Nakayama 引理表明 \(F'\) 在 \(k_s[T]/(F)\) 中可逆。因此 \(F\) 可分，且 \SGARefLink{sga3:VIIB:1.6.A}{1.6.A}~\SGARefLink{sga3:VIIB:1.6.E:item:ii}{(ii)} 表明 \[
  A=k_s[T]/(F)
\] 为满足 \[
  A\widehat\otimes_k\kappa(s)\simeq K.
\] 的 平展 \(k_s\)-代数。由此可知，每个局部预有限 平展 \(k_s\)-代数都是有限秩自由 \(k_s\)-模，因而在 \(k\) 上拓扑自由。所以每个预有限 平展 \(k\)-代数在 \(k\) 上都拓扑自由。

条件 \textup{\SGARefLink{sga3:VIIB:1.6.E:item:v}{(v)}} 蕴含 \SGARefLink{sga3:VIIB:1.6.B}{1.6.B} 的条件 \textup{(d)}，因而蕴含 \textup{\SGARefLink{sga3:vii:paragraph:5-1:equation:i}{(i)}}。这样已证明 \textup{\SGARefLink{sga3:vii:paragraph:5-1:equation:i}{(i)}}、\textup{\SGARefLink{sga3:VIIB:1.6.E:item:iii}{(iii)}} 和 \textup{\SGARefLink{sga3:VIIB:1.6.E:item:v}{(v)}} 等价，并蕴含 \textup{\SGARefLink{sga3:vii:paragraph:5-1:equation:ii}{(ii)}}，后者又蕴含 \textup{\SGARefLink{sga3:VIIB:1.6.E:item:iv}{(iv)}}。

最后，设 \(A\) 为满足 \textup{\SGARefLink{sga3:VIIB:1.6.E:item:iv}{(iv)}} 的预有限 \(k\)-代数。我们证明 \(A\) 在 \(k\) 上形式 平展。可以假设 \(A\) 和 \(k\) 均为局部环，以 \(\kappa\) 表示 \(k\) 的剩余域。由假设，\[
  K=A\widehat\otimes_k\kappa
\] 为 \(\kappa\) 的有限可分扩张，设其次数为 \(n\)。由上述论证及 \SGARefLink{sga3:VIIB:1.6.D}{引理~1.6.D}，存在一个秩为 \(n\) 的自由 \(k\)-代数 \(B\)，它在 \(k\) 上形式 平展，并且有连续态射 \[
  \phi:B\longrightarrow A
\] 使得 \[
  \phi\widehat\otimes_k\kappa
\] 为同构。由于 \(A\) 在 \(k\) 上拓扑平坦，这蕴含 \[
  \operatorname{Coker}(\phi)\widehat\otimes_k\kappa=0
  \qquad\text{且}\qquad
  \operatorname{Ker}(\phi)\widehat\otimes_k\kappa=0.
\]。Nakayama 的 \SGARefLink{sga3:VIIB:0.3}{引理 0.3} 给出 \[
  \operatorname{Coker}(\phi)=0
  =
  \operatorname{Ker}(\phi).
\]。所以 \(\phi\) 是同构 \textup{（参见 \SGARefLink{sga3:vii:proposition:0-2-b}{0.2.B} 的证明）}，这就完成了 \SGARefLink{sga3:VIIB:1.6.E}{1.6.E} 的证明。
